\documentclass[11pt]{article}
\usepackage[margin=1in]{geometry}
\usepackage[T1]{fontenc}
\usepackage{lmodern,amsmath,amssymb,amsthm,mathtools,booktabs,longtable,array}
\usepackage{microtype,xurl,hyperref}
\hypersetup{hidelinks,pdftitle={CM/LM canonical theorem bank}}
\setlength{\parskip}{0.35em}
\setlength{\parindent}{0pt}
\setlength{\emergencystretch}{3em}
\newtheorem{theorem}{Theorem}[section]
\newtheorem{proposition}[theorem]{Proposition}
\newtheorem{lemma}[theorem]{Lemma}
\newtheorem{corollary}[theorem]{Corollary}
\theoremstyle{remark}
\newtheorem{remark}[theorem]{Remark}
\input{notation}
\title{CM/LM Research Consolidation\\Canonical Theorem Bank and Evidence Boundaries}
\author{Research consolidation for Brian Theory}
\date{18 September 2026}
\begin{document}
\maketitle
\begin{abstract}
This bank separates the existing Boolean operator calculus from a chosen
rotation algebra, explicit modal resource contracts, exact query models and
symbolic support interfaces. It consolidates proofs rather than treating
historical computation as authority. The main candidate contributions are a
finite-chain-ring contextuality classification, its leading-layer codebook
synthesis, and exact characteristic-two query bounds. Priority is not certified.
Background imports, finite-only observations and missing empirical evidence are
identified explicitly. Companion registries carry source locations, test scope
and one-primary-paper ownership for 71 claims.
\end{abstract}
\tableofcontents
\newpage
\input{definitions}
\input{model_contracts}

\section{Imported Paper B foundations}
\label{sec:paperb}
\claimid{CM-B-001--007}
\begin{proposition}[Imported CM selection and aligned fusion]
Let $s(x)=(x,1+x)^T$ and let $C_f$ have entries $f(r,t)$ in true-first order.
Then $s(x)^T C_f s(y)=f(x,y)$. If two occurrences have the same declared
ordered operand frame, for every binary Boolean connective $g$,
\[
 g(f_1(x,y),f_2(x,y))=s(x)^T\bigl(g(C_{f_1},C_{f_2})\bigr)s(y),
\]
where $g$ acts entrywise. Operand exchange transposes the CM; input negations
permute the corresponding axis.
\end{proposition}
\prov{S-B sections 2--4 and appendix A.1--A.4. Imported, not re-proved as later
novelty. One-hot selection is the reason arbitrary $g$ is valid; no distributive
law of $g$ over XOR is needed. Structural frame admission, retabulation and
fallback remain different compiler paths.}

\begin{proposition}[Imported symbolic pairing]
The formula-valued LM is
\[
 M_f(X,Y)=\begin{pmatrix}f(X,Y)&f(X,\neg Y)\\
 f(\neg X,Y)&f(\neg X,\neg Y)\end{pmatrix}.
\]
For any formulas $P,Q$, its logical pairing is
\[
 s(P)^T M_f(X,Y)s(Q)\equiv f(P\leftrightarrow X,Y\leftrightarrow Q).
\]
Boolean valuations commute with this pairing. Positive valuation $X=Y=1$
recovers $C_f$.
\end{proposition}
\prov{S-B section 8, Proposition 2 and Theorem 3; appendix A.5. These are
Boolean formulas, not an already selected modal measurement theory. A four-bit
token and an explicit $2^n$-entry truth relation have different cost contracts.}

\section{Rotation lift: classical algebra and its semantic boundary}
\label{sec:rotation}
\claimid{ROT-001, ROT-002, ROT-007}
\begin{theorem}[Cyclic commutant and chain ring]
For the regular $N$-cycle on $\F_2^N$, the centralizer is
$\F_2[R]\cong\F_2[x]/(x^N-1)$. If $N=2^n$ and $u=I+R$, this is
$\F_2[u]/(u^N)$. Multiplication by $u^a v$, for a unit $v$, has binary rank
$N-a$, $0\leq a\leq N$. A coefficient polynomial is a unit iff its coefficient
parity is one. Consequently
\[
 \operatorname{End}_{C_N}((\F_2^N)^m)=\Mat_m(\F_2[R]),\quad
 \operatorname{Aut}_{C_N}((\F_2^N)^m)=\GL_m(\F_2[R]).
\]
\end{theorem}
\begin{proof}
A commuting operator is determined by its first column $v$: its $j$th column
must be $R^jv$. The $N$ powers of $R$ are independent by application to $e_0$,
so the only defining relation is $R^N=I$. Frobenius gives
$(1+x)^N=1+x^N$. In the $u$ basis, multiplication by $u^a$ shifts $N-a$ basis
vectors to independent nonzero vectors and kills the remaining $a$.
Invertibility is equivalent to a nonzero constant $u$ coefficient, namely the
original polynomial evaluated at one. The block commutant statement applies
the same condition to every block.
\end{proof}
\prov{S-ADV D1 and D3. Standard group-algebra/commutant facts, not new CM
algebra. The maximal equivariant gate group is derived; support, tensors and
preparation rules are not.}

\claimid{ROT-003, ROT-005}
\begin{corollary}[Full-degree criterion and distinct filtrations]
For an $n$-variable Boolean function $f$, $n\geq1$, its cyclic coefficient lift
of length $2^n$ is invertible iff $\bigoplus_x f(x)=1$, equivalently the ANF
coefficient of $x_1\cdots x_n$ is one, equivalently $\deg f=n$.
This does not identify the whole cyclic valuation filtration with ANF degree.
\end{corollary}
\begin{proof}
Summing a proper-support monomial over all assignments gives zero in $\F_2$;
the full monomial contributes one. Apply the unit criterion.
For binary-indexed coefficients $a_j$,
\[
 \sum_j a_j(1+u)^j=\sum_k b_k u^k,\qquad
 b_k=\bigoplus_{j\supseteq k}a_j.
\]
The cyclic valuation uses the least \emph{numeric} $k$ with $b_k\ne0$.
By contrast the elementary-abelian group algebra
$\F_2[u_1,\ldots,u_n]/(u_i^2)$ uses total degree, the Hamming weight of $k$.
For an ANF monomial with support $S$, its truth coefficient polynomial there is
$\prod_{i\in S}(1+u_i)\prod_{i\notin S}u_i$; the least degree is $n-|S|$.
Distinct highest-degree ANF monomials give distinct least-degree terms. Thus
its radical order is $n-\deg f$. For a four-cycle the affine coefficient
vectors $(1,1,0,0)$ and $(1,0,1,0)$ have cyclic orders one and two, respectively.
\end{proof}
\prov{S-ADV D1--D2. The Reed--Muller/radical identification is classical
Berman--Charpin theory, not a new finding \cite{Andriatahiny2016}. Cyclic order
also encodes the standard periodic linear-complexity invariant
$N-\deg\gcd(a(x),x^N-1)$. Original Games--Chan priority details remain an
inherited bibliographic lead, not a newly read source.}

\claimid{ROT-004, ROT-006, ROT-009}
\begin{proposition}[Reversal, complement and XOR interferometer]
For $Ke_j=e_{-j}$, $KRK=R^{-1}$ and
$K\Lambda(c)K=\Lambda(c)^T$. In the specified CM cyclic order this reflection
exchanges the two Boolean arguments. Complementing a length-$N=2^n$ truth
vector adds $u^{N-1}$. With
\[
 W=\begin{pmatrix}I&0\\ I&I\end{pmatrix},\qquad
 D_k=\diag(I,R^k),
\]
we have $W^2=I$ and $WD_kW(v,0)^T=(v,v+R^kv)^T$.
\end{proposition}
\begin{proof}
Reflection sends each power $R^j$ to $R^{-j}= (R^j)^T$. Complement adds the
sum of all powers; at length $2^n$ this is $(1+R)^{N-1}$. The displayed
interferometer identity follows by three block multiplications.
\end{proof}
\scope{Complement is translation, not multiplication by the socle. Reversal
is not generally operator inversion, though they coincide on the eight units
of $\A$ at length four. $R$ is unipotent, not multiplication by the complex
number $i$. Its eigenvalues in field extensions are all one; a three-dimensional
unipotent block can already have order four. The ordinary identity CM lifts to
$I+R^2=u^2$, whose square is zero, so this lift cannot preserve the original CM
matrix product or entrywise AND.}
\prov{S-ADV D1--D3; S-CAP F1. Standard identities used as typed motivation.}

\section{Complete-basis contextuality over chain rings}
\label{sec:context}
\claimid{MOD-003}
\begin{lemma}[Residue-hyperplane criterion]
For the complete primitive projective basis support model over a finite
commutative local ring, a global assignment exists iff there are hyperplanes
$H_A\subset k^m$, $H_B\subset k^n$ such that $rMt^T\ne0$ for every primitive
lift $r,t$ of vectors outside $H_A,H_B$, respectively.
\end{lemma}
\begin{proof}
Given a global assignment, let $U_A$ be the rays selected at least once among
Alice's settings. The complement of $U_A$ contains no basis: such a setting
would have to select a ray in its own complement. If the reductions of the
complement spanned $k^m$, choose a residue basis from them; its lifts would be
an invertible local-ring basis, a contradiction. Put their span in a hyperplane
$H_A$. Every lift of a residue ray outside $H_A$ belongs to $U_A$.
Do the same for Bob. Every cross pair from $U_A,U_B$ is selected in some
pair of settings and must have nonzero contraction. Conversely, each invertible
basis has a row reducing outside any fixed hyperplane. Selecting one such row
in each setting on each side gives a compatible assignment by hypothesis.
\end{proof}
\prov{S-ADV D4 Lemma 8. Candidate formulation; priority unresolved. This lemma
uses locality but not the chain condition. The independent dual-number check
uses explicit support and 2-SAT rather than this criterion.}

\claimid{MOD-005}
\begin{lemma}[Valuation-uniform Hardy event]
For $D=\diag(p,q_0)$ with $p=\pi^a$, $q_0=\pi^b\ne0$, $a\leq b$, put
$c=\pi^{b-a}$, so $pc=q_0$. The reversible basis matrices
\[
 F=\begin{pmatrix}0&1\\1&0\end{pmatrix},\quad
 X=\begin{pmatrix}1&0\\1&1\end{pmatrix},\quad
 C_c=\begin{pmatrix}1&0\\-c&1\end{pmatrix}
\]
give, for Alice settings $F,C_c$ and Bob settings $F,X$,
\begin{align*}
 FDF^T&=\begin{pmatrix}q_0&0\\0&p\end{pmatrix},&
 FDX^T&=\begin{pmatrix}0&q_0\\p&p\end{pmatrix},\\
 C_cDX^T&=\begin{pmatrix}p&p\\-q_0&0\end{pmatrix},&
 C_cDF^T&=\begin{pmatrix}0&p\\q_0&-q_0\end{pmatrix}.
\end{align*}
The supported event $(F_A=0,F_B=0)$ has no global extension.
\end{lemma}
\begin{proof}
The determinants are units and direct multiplication gives the matrices.
The seed forces $X_B=1$, then $C_A=0$, then $F_B=1$, contradicting the seed.
The support strings are $1001,0111,1110,0111$.
\end{proof}
\prov{S-ADV D5. The Hardy implication pattern is known; the uniform
valuation-dependent implementation is part of the candidate chain-ring
classification. The minus sign is essential outside characteristic two.}

\claimid{MOD-004}
\begin{theorem}[Smith contextuality trichotomy]
For a nonzero $m\times n$ matrix over $K$, the complete-basis support type is
\[
\begin{array}{c|c}
 r=1 & \text{local}\\
 r\geq2,\ h=1 & \text{logical but not strong}\\
 h\geq2 & \text{strong}.
\end{array}
\]
\end{theorem}
\begin{proof}
Invertible left and right transformations relabel the complete basis families,
so take Smith form. If $r=1$, write $M=\diag(p,0,\ldots)$ and retain the
outcomes of any supported seed event. In all other settings choose a row whose
first coordinate is a unit. New-new pairs are nonzero multiples of $p$.
A seed-new pair is also nonzero: if $pr_1t_1\ne0$, then neither $pr_1$ nor
$pt_1$ is zero, and multiplication by a unit preserves nonzeroness. Thus every
supported seed extends.

If $h=1$, in every basis select a first-coordinate-unit row. Write
$M=\pi^{a_1}\diag(1,\pi^{a_2-a_1},\ldots)$. The contraction in parentheses
is a unit plus radical terms, so the selected cross pair is nonzero. Hence a
global assignment exists.

If $h\geq2$, suppose the hyperplanes in the preceding criterion existed and
write $M=\pi^a M_0$ with $\rank_k\overline M_0=h$. Let $H_B=\ker\beta$.
The vectors outside $H_A$ span $k^m$. Therefore some $x$ outside $H_A$ has
$w=x\overline M_0$ not in the one-dimensional span of $\beta$. There is
$y\in\ker w$ outside $\ker\beta$. Lift $x,y$ to $r,t$. The row $rM_0$ has a
unit coordinate; $rM_0t^T$ lies in the radical. Adjust the corresponding
coordinate of $t$ by a radical element to cancel that dot product exactly.
The residues remain $x,y$, but $rMt^T=0$, contradicting the hyperplane
criterion. Thus no global assignment exists.

Finally, if $r\geq2$, embed the preceding two-coordinate Hardy construction
in the first two nonzero Smith coordinates and extend the settings with
standard rows. The seed and each forcing step remain in those coordinates;
extra outcomes have zero cross entries with the forced rows and cannot evade
the contradiction. Thus some event never extends. The three cases follow.
\end{proof}
\prov{S-ADV D4 Theorem 9; imported and independently checked in S-CAP R1.
Candidate novelty, not certified priority. The proof applies to mixed-characteristic
chain rings as well as polynomial quotients; the symbolic $\F_2$ enrichment later
has a narrower scalar scope.}

\claimid{MOD-006}
\begin{corollary}[Basis counts]
The numbers of primitive projective rows and unordered projective bases in
$K^m$ are
\[
 q^{(s-1)(m-1)}\frac{q^m-1}{q-1},\qquad
 \frac{q^{m^2(s-1)}|\GL_m(q)|}{|K^\times|^m m!},
\]
respectively. For $\A^2$ they are 24 and 192; $|\GL_2(\A)|=24576$.
\end{corollary}
\begin{proof}
There are $q^{sm}-q^{(s-1)m}$ primitive vectors. The unit action is free on
them and $|K^\times|=q^{s-1}(q-1)$. An invertible residue matrix has
$q^{m^2(s-1)}$ lifts, each invertible. Quotient ordered invertible rows by their
independent unit scales and by the free row-permutation action.
\end{proof}
\prov{Standard finite local-ring counting; S-ADV D4.}

\section{Tensor choice, resources and exact protocols}
\label{sec:resources}
\claimid{MOD-001, MOD-002, MOD-009, MOD-010}
\begin{proposition}[Tensor and orbit boundaries]
In $\A\otimes_\A\A$, $u^3\otimes u=0$; in
$\A\otimes_{\F_2}\A$ it is nonzero. Primitive preparations remain primitive
under a balanced tensor, but conditioning the primitive resource
$\diag(1,u)$ by the second standard effect gives the nonprimitive nonzero
vector $(0,u)$. For Smith exponents $(a_i)$,
\[
 \rank_{\F_2}\lambda_d(M)=\sum_i\max(4-a_i,0).
\]
Unrestricted local binary matrix equivalence is classified only by this rank.
\end{proposition}
\begin{proof}
Balanced tensor multiplication identifies the first scalar tensor with
$u^4=0$. A nonzero simple tensor over a field is nonzero, as witnessed by two
nonvanishing dual functionals. A primitive vector has a unit coordinate;
the tensor of two such coordinates is a unit. The stated conditional vector
is obtained by selecting the second matrix row. Regular expansion of Smith
form is block diagonal; each $u^{a_i}$ block has rank $4-a_i$, and invertible
Smith changes remain invertible after expansion. Field matrix equivalence by
independent invertible left/right changes has ordinary rank normal form.
\end{proof}
\scope{$D(0,2)$ and $D(1,1)$ both have binary rank six but different restricted
contextuality. The latter is nonprimitive. Primitive dimension-three examples
$\diag(1,1,u^3)$ and $\diag(1,u,u^2)$ both have binary rank nine but different
leading multiplicities. The particular audited restricted shared/literal
\emph{bipartite support tables} are nevertheless isomorphic: literal Choi
contraction transposes Bob's blocks, and $R\mapsto R^{-1}$ permutes the admitted
basis family. This does not make the tensor theories equivalent.}
\prov{S-ADV D6; S-AUD section 7.4; S-CAP F2. Standard algebra plus a
model-specific restriction comparison.}

\claimid{ALG-001}
\begin{lemma}[Exact linear discrimination]
Let $S_i$ be nonzero state sets in a vector space over a field and
$U_i=\Span(S_i)$. A complete injective linear readout can distinguish the labels
with every possible outcome correct iff $\sum_i U_i$ is a direct sum.
\end{lemma}
\begin{proof}
The coordinates assigned to different labels must be disjoint. The image of
$U_i$ therefore lies in its own coordinate subspace. Injectivity of the collected
readout forbids a nontrivial dependence between these subspaces. Conversely,
concatenate bases of the independent $U_i$ and extend to a basis of the entire
space; the resulting dual-basis readout distinguishes the labels.
\end{proof}
\prov{S-CAP Q1. The single-state independence criterion is modal/linear-algebra
background; see also \cite{DiamondSchumacher2023}. Its uncertainty-set and
recorded-branch form is a useful supporting formulation, not the P02 headline.}

\claimid{RES-001, RES-002}
\begin{theorem}[Leading-layer exact codebook]
Under contract DC, the largest exactly distinguishable invertible-orbit
codebook for $M\in\Mat_d(K)\setminus\{0\}$, $d\geq2$, has $d h$ messages.
Consequently advantage over the $d$-message baseline is equivalent to strong
contextuality under the same complete chain-ring basis contract.
\end{theorem}
\begin{proof}
Write $M=\pi^a N$ with $\rank_k\overline N=h$. The reduction of a divided
codeword $UN$ lies in
$\{X\overline N:X\in\Mat_d(k)\}$, a vector space of dimension $dh$.
Every encoded resource retains a nonzero leading layer. After any invertible
joint decoder, exactly distinguishable codewords have disjoint nonzero supports;
their leading residue vectors are therefore independent. This proves the upper
bound $dh$.

Invertible $d\times d$ matrices over a field span the full matrix algebra.
Off-diagonal units are differences of $I+E_{ij}$ and $I$. For a diagonal unit,
choose a permutation $P$ sending some $j\ne i$ to $i$; the difference of
$P(I+E_{ji})$ and $P$ is $E_{ii}$. Both matrices are invertible.
Thus one can choose $dh$ invertible residue encoders for which
$\overline U_m\overline N$ are independent. Lift them to invertible matrices
over $K$. The columns $\operatorname{vec}(U_mN)$ are independent modulo the
radical, so extend to a basis modulo the radical and lift the remaining columns.
The full lifted matrix is invertible. Its inverse decodes the actual codeword
$\pi^a\operatorname{vec}(U_mN)$ to $\pi^a e_m$, with exactly one possible
outcome. This attains $dh$. The contextuality equivalence follows from the
trichotomy, since $dh>d$ iff $h\geq2$.
\end{proof}
\prov{S-CAP R2 and S-ADV D4. Candidate theorem/synthesis pending priority
clearance. All 14 nonzero $\A$ Smith representatives have explicit independently
checked decoder matrices. This counts messages in a fixed orbit, not entropy.}

\claimid{RES-003, RES-004, RES-007}
\begin{theorem}[Universal exact teleportation and its separation]
Under TP, every nonzero universally correct branch must have invertible
transfer $T_m=M^T E_m^T$. Hence $M$ must be invertible. Conversely, if $M$ is
invertible and a complete basis of invertible reshaped effects is admitted,
correction by $T_m^{-1}$ is exact, including reference correlations.
For $0<a<s$, $\pi^a I_d$ has $d^2$ DC messages and strong contextuality, but
cannot satisfy this teleportation contract.
\end{theorem}
\begin{proof}
If a nonzero branch has nonzero kernel vector $k$, choose $x$ with $T_mx\ne0$.
Then $x$ and $x+k$ are distinct nonzero inputs with the same possible branch
output. No fixed correction can recover both exactly. Thus a nonzero universal
branch is injective. On a finite equal-rank free module an injective square map
is bijective and has a linear inverse. Since $T_m$ factors through $M^T$, its
invertibility implies that of $M$. Completeness and a nonzero resource forbid
all branches being identically zero. In the sufficient case the identity
$T_m^{-1}T_m=I$ proves the claim; tensoring it with a reference identity
preserves correlations. Finally $\pi^a I_d$ has leading multiplicity $d$ but
annihilates nonzero vectors, so apply the codebook theorem and necessity.
\end{proof}
\scope{Exact raw-vector recovery is the stated task, not an unspecified
projective equivalence relation. A singular resource is not rescued by selecting
one purported nonzero universally correct branch. The nilpotent separator is
nonprimitive and depends on allowing that preparation.}
\prov{S-CAP R3; S-ADV D9. The general transfer identity is established modal
background. Its chain-ring resource comparison is the P01 synthesis.}

\begin{proposition}[Explicit admitted analyzers in the supplied carriers]
The four matrices
\[
 \begin{pmatrix}0&1\\1&1\end{pmatrix},\quad
 \begin{pmatrix}1&0\\1&1\end{pmatrix},\quad
 \begin{pmatrix}0&1\\1&0\end{pmatrix},\quad
 \begin{pmatrix}1&0\\0&1\end{pmatrix}
\]
are invertible and their flattened rows form a basis. They give a complete
four-outcome analyzer over $\A$. Their triple field tensor products give a
complete 64-outcome analyzer for a literal eight-dimensional carrier.
\end{proposition}
\begin{proof}
Each determinant is one. If the coefficients of a linear relation are
$\alpha,\beta,\gamma,\delta$, the four entries successively give
$\delta=\beta$, $\gamma=\alpha$, $\beta=0$, $\alpha=0$.
A tensor product of bases is a basis and a tensor product of invertible
matrices is invertible. The correction follows from the preceding theorem.
\end{proof}
\scope{Do not identify a four-outcome logical analyzer with a complete analyzer
for the literal phase register. General-dimensional analyzer sufficiency can
be developed using the invertible-matrix spanning lemma, but the canonical
supplied protocol claim here retains its explicitly admitted carriers.}

\claimid{RES-005, RES-006, RES-008}
\begin{proposition}[Field codebooks and literal extraction]
For a rank-$r$ matrix in $\Mat_d(F)$, $d\geq2$, unrestricted invertible-orbit
encoding and complete field-linear joint readout give exactly $dr$ codewords.
A rank-$r$ resource admits local rectangular filters extracting $I_t$ from $k$
field-tensor copies iff $r^k\geq t$.
\end{proposition}
\begin{proof}
The image of $X\mapsto XM$ has dimension $dr$. Invertible matrices span
$\Mat_d(F)$ by the preceding argument; select a basis from their images and
apply the discrimination lemma. For extraction, reduce $M$ to rank normal
form. Its $k$th tensor has rank $r^k$; select $t$ nonzero coordinate directions
when possible. Conversely rectangular multiplication cannot increase rank.
\end{proof}
\scope{For a regular $2\times2$ $\A$ resource, restricting encoders to
$\GL_2(\A)$ but allowing a fine unrestricted binary decoder gives $2r_{\rm bin}$
codewords (S-CAP R4). This uses the binary span of $\Mat_2(\A)M$ and is an
asymmetric restriction. It can lie below the ordinary eight-letter baseline;
allowing discard/reprepare gives at least $\max(8,2r_{\rm bin})$. Extraction is
an algebraic filter statement, not a distillation probability.}

\section{Setting and form restrictions}
\label{sec:settings}
\claimid{MOD-007}
\begin{theorem}[Two complete binary settings are not strongly contextual]
For any number of parties, a nonzero $\F_2$ tensor with at most two complete
binary basis settings per site has a global assignment. The same holds for
two-dimensional local modules over any finite commutative local ring with
residue field $\F_2$, using nonzero-amplitude support.
\end{theorem}
\begin{proof}
Two distinct binary dual bases share a row $a$; write the other rows as $b$
and $a+b$. Expand the tensor in coordinates dual to each pair $(a,b)$.
Choose an inclusion-minimal set $S$ of $b$ coordinates having nonzero coefficient.
Outside $S$, select $a$ in both settings. Inside $S$, select $b$ in the first
and $a+b$ in the second. Each joint choice evaluates to the coefficient for
$S$ plus coefficients of proper subsets of $S$. By minimality the latter are
zero and the former is one. Every chosen event is therefore supported.
Coincident settings can be duplicated and require no additional choice.

For a finite local ring with radical $J$, take a nonzero leading layer of the
tensor in $J^a/J^{a+1}$. Apply an $\F_2$-linear functional on that layer which
makes the resulting coefficient tensor nonzero. Reduce the local bases modulo
$J$ and apply the field construction. Each selected contraction then has a
nonzero leading-layer functional, so its original ring coefficient is nonzero.
\end{proof}
\scope{Literal independent phase blocks can be organized in the local algebra
$\F_2[u_1,\ldots,u_n]/(u_i^4)$; the same leading-layer proof applies to the
specified coarse measurements without identifying it with shared $\A$.
This theorem permits two-setting Hardy logical contextuality and does not cover
three settings, higher outcomes or arbitrary residue fields.}
\prov{S-CAP R1b. Candidate general theorem. The source algorithm checked
65,808 nonzero states for one through four mobits and 1,050,666 supported-context
checks; the proof, not that enumeration, handles arbitrary party count.}

\claimid{RES-009, ROT-008}
\begin{proposition}[Transpose-orthogonal span and a separate form obstruction]
For $d\geq2$,
\[
 \Span_{\F_2}O_d(\F_2)=
 \{M:M\mathbf1=c\mathbf1,\ \mathbf1^TM=c\mathbf1^T
 \text{ for some }c\in\F_2\},
\]
of dimension $(d-1)^2+1$. Therefore the stated complete rank-one teleportation
analyzer with transpose-orthogonal corrections cannot exist. Separately, the
native $W$ and $\diag(I,R)$, $R\ne I$, preserve no common nondegenerate
bilinear form on the same doubled space.
\end{proposition}
\begin{proof}
Every column of an orthogonal binary matrix has odd weight; therefore both its
row and column sums are one. Linear combinations lie in the displayed space.
Permutation matrices are orthogonal. Differences of permutations produce each
four-corner rectangle involving a fixed last row and column; these span the
$(d-1)^2$-dimensional zero-row/column-sum space. Adjoining $I$ adds the common-sum
direction. For a fixed invertible teleportation resource, reshaped effects with
orthogonal corrections lie in an invertible translate of this span, too small
to span the $d^2$-dimensional dual space.

For the second statement write a form matrix as
$\begin{pmatrix}A&C\\E&F\end{pmatrix}$. Invariance under $W$ gives $F=0,E=C$.
Nondegeneracy forces $C$ invertible, since $(0,x)$ is a null vector when
$Cx=0$. Invariance under $\diag(I,R)$ gives $CR=C$, hence $R=I$, a contradiction.
\end{proof}
\prov{S-ADV D7--D8 and S-CAP E1--E2. The permutation-span ingredient is
classical \cite{Lueneburg1988}; the dimension-eight span is 50, not merely
bounded by 63. These are scoped model obstructions, not universal claims about
all symplectic or finite-field theories. $R$ itself preserves the alternating
form $R^2$; different admissible gate families remain possible.}

\section{Characteristic-two exact query theory}
\label{sec:query}
\claimid{ALG-002}
\begin{theorem}[Deutsch: a lower bound and its properly typed upper bound]
With any fixed invertible pointwise $G_0,G_1$ over a characteristic-two field,
no exact one-query linear modal procedure computes $f(0)+f(1)$. For the standard
XOR oracle QXOR the exact query complexity is two.
\end{theorem}
\begin{proof}
Write the prepared state by address sectors $(a_0,a_1)$, allowing all target and
ancillary degrees of freedom. The four postquery states satisfy
\[
 v_{00}+v_{11}=v_{01}+v_{10}
 =((G_0+G_1)a_0,(G_0+G_1)a_1)=w.
\]
If $w\ne0$, the constant and balanced spans intersect nontrivially, contrary
to exact discrimination. If $w=0$, both sector differences vanish and all four
states are identical. Complete subsequent linear readout cannot remove either
obstruction. Under QXOR, query both addresses on basis labels, retain the two
answers and compute their XOR; this attains two.
\end{proof}
\scope{The upper bound is not valid for arbitrary $G_0,G_1$ without an
informative interface. If $G_0=G_1$, every oracle is identical at every query
count. This explicit qualification is a consolidation correction to the broad
wording of S-CAP Q2.}
\prov{S-CAP Q2, read with the QXOR contract; ALG-002; correction C17. Candidate
exact-query result after dedicated but nonexhaustive prior-art search.}

\claimid{ALG-003}
\begin{theorem}[Deutsch--Jozsa one-query obstruction]
For $n\geq2$ and the constant-or-balanced promise on $N=2^n$ addresses, the
same pointwise characteristic-two contract has no exact one-query solution.
For QXOR the established bounds are $2\leq Q_{\rm exact}\leq N/2+1$.
\end{theorem}
\begin{proof}
Partition the addresses into four equal blocks $A,B,C,D$. The balanced
functions with supports $A\cup B$, $A\cup C$, $B\cup C$ have pointwise XOR
zero. The pointwise oracle is affine in its control bit, so their three oracle
operators sum to $O_0$ in characteristic two. The nonzero constant-zero
postquery state is therefore in the balanced-state span, violating exact
discrimination. The QXOR upper bound queries $N/2+1$ distinct addresses: agreement
forces constancy, and any difference proves balance. The $n=1$ case is Deutsch.
\end{proof}
\scope{No claim is made that the classical upper bound is optimal for all
multi-query modal procedures. One-query failure is stronger than failure of
a particular Hadamard circuit, but is not an all-algorithm lower bound of
$N/2+1$.}
\prov{S-CAP Q3. Candidate query result; source checks cover 56 operator
coordinates, not every ancillary state. The proof supplies ancillary generality.}

\claimid{ALG-004}
\begin{theorem}[Exact Bernstein--Vazirani lower bound]
For $f_s(x)=s\cdot x$, $s\in\F_2^n$, a $q$-query exact characteristic-two
linear modal algorithm under the fixed pointwise oracle requires $q\geq n$.
Under QXOR, $n$ basis-label queries attain the bound.
\end{theorem}
\begin{proof}
The oracle is affine in the secret:
$O_s=O_0+\sum_j s_j A_j$. After $q$ queries, every state component is a
multilinear polynomial of degree at most $q$ in the bits of $s$ (replace
$s_j^2$ by $s_j$). Thus all $2^n$ final states lie in the span of at most
\[
 \sum_{j=0}^{\min(q,n)}\binom nj
\]
coefficient vectors. For $q<n$ this is strictly less than $2^n$, while exact
separation of the $2^n$ singleton labels requires independent final states.
Ancillas enlarge coefficient vectors, not their number. Recording all complete
adaptive branches in a direct sum preserves both the degree bound and the
injectivity requirement. QXOR queries at $e_j$ directly return $s_j$.
\end{proof}
\prov{S-CAP Q4. Polynomial-method architecture is established
\cite{Beals2001}; the specific modal exact bound needs priority clearance.
In complex-amplitude computation, XOR of secret bits is not represented by this
same degree-one scalar polynomial, so there is no contradiction with ordinary
one-query BV.}

\claimid{ALG-005, ALG-006, ALG-007}
\begin{proposition}[Kickback and the missing character basis]
In $\A$, set $z=R^2=1+u^2$. For the standard target flip $X$ and
$\chi=(1,z)^T$, $X\chi=z\chi$. But
\[
 H_z=\begin{pmatrix}1&1\\1&z\end{pmatrix}\sim\diag(1,u^2)
\]
is not invertible. Its $n$th tensor over $\A$ has binary rank $4+2n$.
Multiplicative characters of $C_2^n$ into a characteristic-two field are all
trivial; an $\A$-valued character matrix reduces modulo $u$ to all ones.
\end{proposition}
\begin{proof}
Use $z^2=1$. Row and column elimination give the Smith diagonal displayed.
Tensoring the Smith forms yields $u^{2j}$ with multiplicity $\binom nj$;
only $j=0,1$ survive, with binary ranks four and two. In a characteristic-two
field $a^2=1$ implies $(a+1)^2=0$ and therefore $a=1$. Every group element
in $C_2^n$ has order at most two, proving character triviality. Every unit of
$\A$ reduces to one, proving singular reduction for the ring character matrix.
\end{proof}
\scope{The bit-flip target cannot have a primitive eigenvector of eigenvalue
$R$, because applying the flip twice would give $u^2\chi=0$. A four-level cyclic
target has eigenvector entries $R^{-j}$ and can kick back $R$, but its oracle is
a distinct interface. Compute-control-uncompute generally costs two ordinary
bit queries. Literal independent registers do not move an $\A$ scalar between
factors automatically. The no-go is not about all finite-field transforms or
odd-order groups.}
\prov{S-CAP Q5. Standard modular-character obstruction plus model-specific
rank/algorithm comparison.}

\claimid{ALG-008, ALG-009, ALG-010, ALG-011, ALG-012}
\paragraph{Finite and supporting algorithm records.}
The Simon record exhausts 65,535 nonzero preparations in dimension 16 for a
specified one-query, no-ancilla, $n=2$ contract with 36 standard oracles and
three shifts. It is not an arbitrary-size/adaptive Simon theorem. The
conventional Grover sign and diffusion construction collapses in characteristic
two, and support alone supplies no probability to amplify; this is not an
all-search impossibility proof. The known Willcock--Sabry UNIQUE-SAT protocol
\cite{WillcockSabry2011} is a one-query positive control on the none-or-one-marked
promise with $O(n)$ other gates; 132 promised instances for $n=1,\ldots,6$ were
reproduced. Nonlinear Boolean $f$ still embeds as a reversible basis permutation
$\lvert x,y\rangle\mapsto\lvert x,y+f(x)\rangle$; linear amplitude evolution
does not force $f$ to be linear. Compute-copy-uncompute is standard reversible
programming \cite{Bennett1973,James2011}. Printing a full state description or
support table is not an allowed cost-free operational readout in QXOR.

\section{The enriched LM interface}
\label{sec:bridge}
\claimid{LMB-001, LMB-002}
\begin{proposition}[Two obstructions to a bare identification]
The fixed polarity-orbit LM family is not closed under ordinary Boolean matrix
composition. A unital Boolean-ring homomorphism to a local ring has image
contained in $\{0,1\}$ and cannot produce nontrivial nilpotent amplitudes.
\end{proposition}
\begin{proof}
Take $M=\lvert X\rangle\langle Y\rvert$. Then
\[
 M^2=\begin{pmatrix}XY&0\\0&\neg X\neg Y\end{pmatrix},
\]
which is not the polarity orbit of its upper-left entry. The finite checker
finds 144 of 256 products outside the fixed family. For the second statement,
the image of every Boolean element is an idempotent. In a local ring, if
$e^2=e$, one of $e,1-e$ is a unit; $e(1-e)=0$ then forces $e=1$ or $e=0$.
\end{proof}
\prov{S-LMB, independently checked here. Local-ring idempotents are standard;
the conclusion rules out a particular canonical identification, not useful
coefficient encodings.}

\claimid{LMB-003, LMB-004, LMB-005, LMB-007, LMB-008}
\begin{theorem}[Sufficient symbolic support bridge]
After choosing the $\F_2$-algebra $\A$, let
$\Bool_\A=\Bool\otimes_{\F_2}\A\cong\Bool[u]/(u^4)$.
Every valuation $v:\Bool\to\F_2$ extends as
\[
 v_\A\left(\sum_{j=0}^3 f_j u^j\right)
 =\sum_{j=0}^3v(f_j)u^j.
\]
It commutes with enriched matrix and tensor contractions. For
$q=\sum_j f_j u^j$ put $\Supp(q)=\bigvee_j f_j$. Then for any chosen
compatible enriched effect $e$ and state $\psi$,
\[
 v\bigl(\Supp(e\psi)\bigr)=1
 \quad\Longleftrightarrow\quad v_\A(e)v_\A(\psi)\ne0.
\]
\end{theorem}
\begin{proof}
The $\F_2$ basis $1,u,u^2,u^3$ gives unique enriched coefficients. Base change
extends $v$ uniquely as an $\A$-algebra map and preserves finite products and
sums, hence contractions. A valued $\A$ coefficient is nonzero exactly when at
least one of its four binary coefficients is one. A Boolean valuation preserves
OR, giving the equivalence.
\end{proof}
\scope{The theorem is conditional on choosing the scalar algebra and effects;
``nonzero means possible'' is an added interpretation. Bare valued logical bras
yield only the two one-hot rows, not the complete $\A$ effect family. The
nonzero map is not additive ($a+a=0$) and not multiplicative in the presence of
zero divisors. For a finite valuation set $\Omega$, this enrichment is the
extensional power $\A^\Omega$, not an efficient compression theorem. General
mixed-characteristic $K$ need not admit this $\F_2$ tensor construction.}
\prov{S-LMB; S-CAP F8; ordinary base change and homomorphism preservation.
Useful semantic interface, not a derivation of all measurement postulates.}

\claimid{LMB-006, NEG-002}
\begin{proposition}[Division-free basis branches]
For an invertible coordinate readout $B_0$ over $\Bool_\A$ and coordinate
projectors $P_i$, define $J_i=B_0^{-1}P_iB_0$. Then
\[
 J_i^2=J_i,\quad J_iJ_j=0\ (i\ne j),\quad\sum_iJ_i=I.
\]
The branch $J_i\psi$ is nonzero iff the $i$th readout coordinate is nonzero;
row-unit rescaling leaves $J_i$ unchanged. Valuation commutes with the branch
formula whenever $B_0$ is symbolically invertible.
\end{proposition}
\begin{proof}
Conjugate the coordinate-projector identities. A column of $B_0^{-1}$ is
unimodular, since some row of $B_0$ pairs with it to one. Multiplying that column
by a nonzero scalar cannot give zero, proving branch/nonzero equivalence.
Diagonal row-unit scaling commutes with $P_i$ and cancels in the conjugation.
Valuation preserves the inverse identity and therefore the full expression.
\end{proof}
\scope{No division by an outcome amplitude is used. No probability,
normalization or closure of primitive preparations under conditioning follows.
A category of instruments or effectus requires additional axioms and proofs.}

\section{Further supporting no-go and finite evidence records}
\label{sec:diagnostics}
\claimid{RES-010, RES-011, RES-012}
\begin{proposition}[Cloning, qualified deletion and exact syndromes]
For $V=\F_2^d$, $d\geq2$, no linear map clones every nonzero vector as
$v\mapsto v\otimes v$. No injective deterministic linear map deletes a copy
as $v\otimes v\mapsto v\otimes b$ with fixed nonzero blank $b$ and no
input-dependent environment. For an injective code map $E$ and errors $F_a$,
an exact decoder retaining distinct syndromes exists on their joint image iff
each $F_aE$ is injective and their images form a direct sum.
\end{proposition}
\begin{proof}
Cloning $e_i+e_j$ requires cross terms absent from the linear sum of the two
clones. Duplicated vectors span the subspace generated by $e_i\otimes e_i$
and $e_i\otimes e_j+e_j\otimes e_i$, of dimension $d(d+1)/2>d$; their proposed
deleted outputs span at most $d$ dimensions, contradicting injectivity.
For syndromes, distinct retained labels occupy independent subspaces and must
preserve each input, proving necessity. Conversely use the inverse of each
$F_aE$ on its own independent image and append its label; these maps combine
to an isomorphism on the direct sum.
\end{proof}
\scope{Linearity alone does not forbid deletion: the noninjective map
$D(e_i\otimes e_j)=\delta_{ij}e_i\otimes b$ deletes all duplicated $\F_2$
vectors. It is not automatically a complete deterministic allowed process.
The syndrome criterion is not the general degenerate quantum error-correction
criterion. Remote preparation and swapping are the standard contraction
identities $\im M^T$ and $MEN$ in compatible conventions; over rings,
primitive preimage conditions still matter.}
\prov{S-CAP F3--F6. Generic modal no-go results are prior art
\cite{Schumacher2010,DiamondSchumacher2023}; syndrome distinctions are
supporting algebra, not a fault-tolerance threshold claim.}

\claimid{MOD-008, RES-013, NEG-001, NEG-003, NEG-004}
\paragraph{Finite GHZ and probability records.}
The tested three-setting GHZ support has 27 contexts, 112 possible events and
no compatible assignment among 512 assignments. Within that family, 101,583
subsets of at most five contexts were checked without contradiction; a
six-context witness is $ZZZ,ZXX,XZY,XXX,YYZ,YYY$. This is a finite diagnostic,
not a theorem that every GHZ scenario requires six contexts. The restricted
orthogonal-basis version has global assignments.

For the specified nine-context Bell support, 24 events are possible, but
normalized nonnegative real no-signalling constraints force six of them to
zero. The supplied exact affine constraints are the certificate; the LP is a
check. A weak completion with 18 positive events and a CHSH value of four is
not a faithful probability model of all original events, still less a Born
rule. Schumacher--Westmoreland already discuss weak versus strong probability
resolutions \cite{Schumacher2012}; this is retained as background validation,
not new foundational priority. Historical signed $C_8/C_{16}$ claims are not
silently imported into the Boolean core.

\section{Compiler evidence is not an algebraic performance theorem}
\label{sec:compiler}
\claimid{CMP-001--006, NEG-005, NEG-006}
The supplied raw file has 10,500 rows, with 10,389 fully admitted four-arm
comparisons and 111 other status rows. Every fully admitted case ties CM and
generic two-support optimization on final degree, final monomials,
multiplications, output unique nodes, pair terms, peak monomials and serialized
bytes. The independent recount found no differences. S1's SC-A symbolic-work
mechanism gate passes; the corrected S2 construction verifies requested support
rather than retaining the superseded generator's 55.161 percent rate.

These observations show that avoiding early ANF expansion can save work, but
not that the final simplification is uniquely accessible to CM. Equal reductions
do not imply equal recognition overhead. The source package does not supply a
completed natural-workload SC-B/C/D campaign, a confirmatory SC-E abstention
result, or Paper B's separate frozen Gate P evaluation. No fresh timing campaign
was run in this consolidation. The 12 benchmark tests are distinct from Paper
B's reported focused 105-test implementation suite, which was not rerun here.

\section{Reproducibility and remaining proof/priority boundaries}
An independent standard-library checker authored for this consolidation
reconstructs all 65,536 $\A$ resource ranks and Smith classes, classifies all
255 nonzero dual-number resources by explicit support plus SCC 2-SAT, checks
14 actual dense decoders and enumerates all 256 fixed-LM products. The dual-number
counts are 81 local, 72 logical-not-strong and 102 strong, using 62,463 SAT calls.
This is a different-method finite check of the generalized theorem.

Fresh named pytest totals are 31 audit, 59 across four historical Boolean suites,
and 12 benchmark tests, hence 102. Other source scripts perform exact scans and
are recorded separately. An interrupted aggregate driver is retained as such;
individual successful stages, not the interrupted driver, support the reported
results. Original input hashes, stage logs and explicit source locations are
provided in the companion repository. Presence of a file copied from a rerun
working directory alone does not certify that its producing stage was rerun.

No novel theorem is claimed merely because targeted web searches found no exact
match. The finite-chain-ring classification and codebook synthesis, two-setting
obstruction and exact characteristic-two query bounds remain the principal
candidates for specialist priority review. Drafts can be built from these
proofs now, but submission-strength priority and compiler performance language
must respect the remaining gaps.

% Static fallback is generated from the same records as references.bib.
% For a journal BibTeX build, replace this conditional with the two fallback lines.
\IfFileExists{../literature/references_static.tex}{\input{../literature/references_static}}{
\bibliographystyle{plain}
\bibliography{../literature/references}}
\end{document}
